% Part: methods
% Chapter: proofs
% Section: cant-do-it

\documentclass[../../../include/open-logic-section]{subfiles}

\begin{document}

\olfileid[pt]{mth}{prf}{cnt}

\olsection{Não Consigo Fazer Isso!{}}

Todos nós chegamos a um ponto em que sentimos vontade de desistir. Mas
você \emph{consegue} fazer isso. Seu professor e o monitor, bem como
seus colegas, podem ajudar. Peça ajuda a eles!{}  Aqui estão algumas
dicas para ajudá-lo a evitar uma crise, e o que fazer se você sentir
vontade de desistir.

Para garantir que você consiga resolver os problemas com sucesso, faça
o seguinte:
\begin{enumerate}
\item \emph{Comece com a maior antecedência possível.} Ficamos
  ocupados ao longo do semestre e muitos de nós lutam com a
  procrastinação; uma das melhores coisas que você pode fazer é começar
  as tarefas de casa cedo. Dessa forma, se você travar, terá tempo
  para procurar uma solução (que não seja chorar).
\item \emph{Converse com seus colegas de turma}.  Você não está
  sozinho. Outros na turma também podem ter dificuldades---mas podem
  ter dificuldades com coisas diferentes. Conversar sobre o assunto com
  seus colegas pode lhe dar uma perspectiva diferente sobre o problema
  que pode levar a um avanço.  É claro, não copie simplesmente a solução
  deles: peça uma dica, ou explique onde você travou e peça o próximo
  passo. E quando você conseguir, retribua. Ajudar outra pessoa, e
  explicar as coisas, ajudará você a entender melhor também.
\item \emph{Peça ajuda.} Você tem muitos recursos à sua
  disposição---seu professor e o monitor estão lá por você e
  \emph{querem} que você tenha sucesso. Eles devem ser capazes de
  ajudá-lo a resolver um problema e identificar em que ponto do processo
  você está com dificuldade.
\item \emph{Faça uma pausa.} Se você travar, \emph{pode} ser porque
  ficou olhando para o problema por tempo demais. Faça uma pausa curta,
  tome uma xícara de chá, ou trabalhe em um problema diferente por um
  tempo, depois volte ao problema com a mente renovada. Durma sobre o
  assunto.
\end{enumerate}

Percebe como essas estratégias exigem que você tenha começado a
trabalhar na prova com bastante antecedência? Se você começou a prova
às 2 da madrugada do dia anterior à entrega, elas talvez não sejam tão
úteis.

Isso pode soar como pessimismo e desgraça, mas resolver uma prova é um
desafio que compensa no final. Algumas pessoas fazem isso como
profissão---portanto deve haver algo de prazeroso nisso. Como
praticamente tudo, resolver problemas e fazer provas é algo que exige
prática.  Você pode ver colegas que acham isso fácil: provavelmente
eles apenas já tiveram muita prática.  Tente não ceder com facilidade
demais.

Se você ficar mesmo sem tempo (ou sem paciência) em um problema
específico: tudo bem. Isso não significa que você seja burro ou que
nunca vá conseguir. Descubra (com seu professor ou outro estudante) como
se faz, e identifique onde você errou ou travou, para poder evitar fazer
isso da próxima vez que encontrar uma questão semelhante.  Depois tente
fazê-lo sem olhar a solução.  E na próxima vez, comece (e peça ajuda)
mais cedo.

\end{document}
